% !Mode:: "TeX:UTF-8"
\begin{conclusions}

本文围绕"中性原子量子计算体系的光子—原子量子接口"这一方向，建立了端到端第一性原理仿真框架，主要工作与结论如下：

（1）建立了面向中性原子处理节点的端到端量子线路化模型，将发射、量子频率转换、光纤偏振演化、分束干涉与探测等模块统一纳入时间 bin 离散的量子演化描述，实现了从器件参数到远程纠缠性能的定量映射。

（2）实现了基于矩阵积态（MPS）与时间演化块消去（TEBD）的数值求解框架，通过量子轨迹采样处理测量条件化，在可控计算资源下完成数百时间 bin 规模的仿真，输出成功率、条件态保真度等可直接对照实验的统计量。

（3）系统分析了 QFC 噪声、光纤损耗、偏振漂移、探测效率等关键参数对接口性能的影响，给出了成功率—保真度的联合依赖关系与误差来源分解，为接口工程优化提供了定量依据。

本文工作的局限性在于：当前模型采用了若干简化假设（如忽略滤波器件的时间记忆效应、PMD 的频率相关性等），数值方法在高纠缠场景下的键维度增长仍需进一步优化。

未来研究方向包括：（1）多路复用与多节点网络的扩展建模；（2）与量子纠错协议耦合的资源评估；（3）面向片上集成与高效率腔/波导接口的模型推广。

\end{conclusions}
